\chapter*{Resumen}
\addcontentsline{toc}{chapter}{Resumen}

La previsión de demanda mejora cuando el modelo se entrena con más datos, pero los operadores de
distribución que competirían por esos datos no pueden compartirlos. El aprendizaje federado propone
una salida: entrenar un modelo conjunto intercambiando parámetros en lugar de datos. Este trabajo
evalúa si esa promesa se sostiene empíricamente en la previsión de ventas de un supermercado.

Sobre un histórico real de 54 tiendas y 33 familias de producto, agregado a nivel semanal, se
comparan cinco condiciones —modelo local por tienda, centralizado por silo, centralizado global,
federado y federado con personalización— entre sí y frente a cuatro métodos convencionales de
previsión, con idéntica partición temporal, idéntico protocolo de evaluación y contraste estadístico
de todas las diferencias mediante el test de Wilcoxon pareado con corrección por comparaciones
múltiples.

Los resultados contradicen la hipótesis de partida en dos puntos. El modelo centralizado, que la
literatura suele tomar como techo de rendimiento, resulta ser el peor de los cinco: la precisión
empeora de forma monótona a medida que se agregan datos de tiendas heterogéneas en un modelo de
capacidad limitada. Y el modelo federado, por sí solo, es significativamente peor que el modelo
local; es la personalización posterior por tienda la que lo convierte en la mejor de las cinco
condiciones, con un WMAPE de 0,1379 frente a 0,1476 del modelo local. Frente a los métodos no
federados, el balance es más matizado: la mejor condición federada supera por muy poco a una media
móvil de cuatro semanas y queda un 3,95\,\% por detrás de un modelo LightGBM entrenado sobre datos
agrupados, diferencia que el contraste confirma como real.

El sistema se desplegó además sobre tres máquinas virtuales independientes en la nube, con
aislamiento efectivo de los datos de cada silo, verificando que el comportamiento observado en
simulación se mantiene sobre infraestructura distribuida.

En el plano económico, la colaboración federada reduce un 16,50\,\% el error de previsión ponderado
por volumen respecto a que cada tienda entrene en solitario, lo que se traduce en un ahorro anual
estimado de entre 0,5 y 3,9 millones de euros para la red analizada. El beneficio, sin embargo, se
concentra en las series de mayor volumen y casi la mitad de las series empeoran, lo que aconseja
seleccionar el modelo serie a serie. Se argumenta finalmente que, para la tarea concreta de
entrenar un modelo predictivo entre partes que no ceden datos, el aprendizaje federado es un
primitivo más adecuado que un \emph{data clean room}, sin presentarlos como alternativas
excluyentes.

\vspace{1em}
\noindent\textbf{Palabras clave:} aprendizaje federado, previsión de demanda, distribución
comercial, datos no IID, personalización de modelos, MLOps.

\chapter*{Abstract}
\addcontentsline{toc}{chapter}{Abstract}

Demand forecasting improves when models are trained on more data, yet retailers who compete for
that data cannot share it. Federated learning offers a way out: training a joint model by
exchanging parameters instead of raw data. This work evaluates whether that promise holds
empirically for supermarket sales forecasting.

Using a real sales history of 54 stores and 33 product families aggregated weekly, five conditions
—per-store local model, per-silo centralised, globally centralised, federated, and federated with
personalisation— are compared against each other and against four conventional forecasting methods,
under an identical temporal split, an identical evaluation protocol, and paired Wilcoxon testing
with correction for multiple comparisons.

The results contradict the initial hypothesis on two counts. The centralised model, usually treated
in the literature as the performance ceiling, turns out to be the worst of the five: accuracy
degrades monotonically as heterogeneous store data are pooled into a capacity-limited model. And
the federated model on its own is significantly worse than the local one; it is the subsequent
per-store personalisation that makes it the best of the five conditions, reaching a WMAPE of 0.1379
against 0.1476 for the local model. Against non-federated methods the picture is more nuanced: the
best federated condition only marginally outperforms a four-week moving average and falls 3.95\,\%
short of a LightGBM model trained on pooled data, a gap the statistical test confirms as real.

The system was additionally deployed across three independent cloud virtual machines with effective
per-silo data isolation, confirming that the behaviour observed in simulation carries over to
distributed infrastructure.

Economically, federated collaboration reduces volume-weighted forecast error by 16.50\,\% relative
to each store training alone, which translates into an estimated annual saving of between 0.5 and
3.9 million euros for the network studied. The benefit is, however, concentrated in high-volume
series and almost half of the series get worse, which argues for selecting the model series by
series. Finally, it is argued that for the specific task of training a predictive model across
parties that do not share data, federated learning is a more suitable primitive than a data clean
room, without presenting the two as mutually exclusive.

\vspace{1em}
\noindent\textbf{Keywords:} federated learning, demand forecasting, retail, non-IID data, model
personalisation, MLOps.
